\clearpage
\section{MC core centers and shifted primary windows}
\label{sec:corecenter}
This update, dated 23 September 2026, centers the observed fit and the
excluded GP training interval on the reconstructed MC core, rather than
the generated pole. The reconstructed template itself is unchanged.
For a generated hypothesis $m_0$ and an MC-only core location $c(m_0)$,
\begin{equation}
 I(m_0)=[c(m_0)-2.25\sigma_{\rm nominal}(m_0),\,
         c(m_0)+2.25\sigma_{\rm nominal}(m_0)].
\end{equation}
The half-width retains the original nominal resolution. The fitted core
width below does not replace it. No additional guard is used.

The center is determined without observed counts. A Gaussian smoothing
of width $0.2\sigma_{\rm nominal}$ locates the largest MC peak within
$m_0\pm3\sigma_{\rm nominal}$. A bin-integrated Gaussian plus a
nonnegative affine pedestal is then fitted by Poisson deviance to the
unsmoothed template within $\pm1.5\sigma_{\rm nominal}$ of that mode.
The Gaussian mean defines $c$; the pedestal limits the influence of a
local broad component. The mean may vary by one nominal sigma around
the mode, and the fitted width by 0.25--2.5 nominal sigma. All 201 fits
converge without these bounds being active. This is an operational core
definition, not proof of a pure resonance response.

\begin{table}[H]\centering\scriptsize
\begin{tabular}{rrrrrrr}
\toprule
$m_0$ & $c(m_0)$ & $c-m_0$ & $\sigma_{\rm core}$ & MC spread & Max. $|\Delta c|$ & Max. UL change (\%) \\
\midrule
60 & 58.792 & -1.208 & 1.550 & 0.034 & 0.049 & 0.00 \\
80 & 78.103 & -1.897 & 1.637 & 0.007 & 0.003 & 0.00 \\
100 & 97.688 & -2.312 & 1.853 & 0.006 & 0.027 & 0.00 \\
120 & 117.319 & -2.681 & 2.112 & 0.003 & 0.039 & 0.00 \\
140 & 136.956 & -3.044 & 2.392 & 0.006 & 0.052 & 0.29 \\
160 & 156.658 & -3.342 & 2.713 & 0.007 & 0.081 & 0.00 \\
180 & 176.393 & -3.607 & 3.120 & 0.008 & 0.096 & 0.00 \\
200 & 196.182 & -3.818 & 3.577 & 0.011 & 0.130 & 0.00 \\
220 & 215.941 & -4.059 & 4.233 & 0.012 & 0.162 & 4.01 \\
240 & 235.617 & -4.383 & 4.981 & 0.022 & 0.132 & 3.75 \\
\bottomrule
\end{tabular}

\caption{Native-MC core centers and stability. All mass quantities are
in MeV. MC spread is the standard deviation of the estimated center
under 32 independent-bin Poisson resamples. The last two columns vary
the core-fit half-width from its default 1.5 nominal sigma to 1.25 and
2.0, and refit the observed MC limit using each alternative center.}
\end{table}

The maximum center-definition displacement is 0.162~MeV and the largest
corresponding limit change is 4.01\%. Native-bin masks are discrete;
many smaller shifts select exactly the same bins. Bin resampling
quantifies conditional center precision only; it does not establish
candidate independence or cover production systematics.

The direct-template centers shift by 1.21--4.38~MeV below their generated poles. The median absolute MC/Gaussian limit difference in the matched centered windows is 17.83\%, compared with 22.05\% for the original paired prescription. The native new/old MC limit ratios span 0.822--1.201. The smallest native MC local probability is at $m_0=80$~MeV, $c=78.103$~MeV: $p_0=0.00787$ ($Z_0=2.415$), compared with $p_0=0.02477$ before shifting the window.


\clearpage
\subsection{Observed scans after centering}
The GP mean and covariance are recomputed from sidebands outside each
shifted interval. Archived kernel coordinates remain anchored at $m_0$.
The inherited density conversion, nominal resolution and branching
conversion also retain $m_0$, isolating the center change from a change
of coupling convention. Both the MC and its new Gaussian benchmark use
the same shifted window and GP constraint. The benchmark Gaussian is
centered at $c(m_0)$ with width $\sigma_{\rm nominal}(m_0)$; it is not
the historical Gaussian fit evaluated at mass $c$ with new kernel or
density settings. All curves are plotted against generated hypothesis
$m_0$, not reconstructed center $c$.

\fig{.71}{../figures/core_centered_comparison.pdf}{Original and
MC-centered observed inference for the 2021 10\% sample. Native MC
points are marked; continuous MC curves include exploratory morphs.
The bottom panels compare only the ten generated masses. The matched
Gaussian comparison separates the remaining shape difference from
placing a Gaussian peak at the wrong reconstructed position.}

\clearpage
\subsection{Native-mass limits and local probabilities}
\begin{table}[H]\centering\scriptsize
\begin{tabular}{rrrrrrr}
\toprule
$m_0$ & Old MC limit & New MC limit & New/old & Old MC $p_0$ & New MC $p_0$ & New G $p_0$ \\
\midrule
60 & $1.41\times10^{-5}$ & $1.16\times10^{-5}$ & 0.822 & 0.49502 & 0.50000 & 0.50000 \\
80 & $6.36\times10^{-6}$ & $7.64\times10^{-6}$ & 1.201 & 0.02477 & 0.00787 & 0.00348 \\
100 & $2.09\times10^{-6}$ & $2.30\times10^{-6}$ & 1.103 & 0.50000 & 0.50000 & 0.50000 \\
120 & $2.26\times10^{-6}$ & $2.30\times10^{-6}$ & 1.021 & 0.50000 & 0.50000 & 0.50000 \\
140 & $1.96\times10^{-6}$ & $2.09\times10^{-6}$ & 1.067 & 0.50000 & 0.50000 & 0.50000 \\
160 & $4.79\times10^{-6}$ & $5.04\times10^{-6}$ & 1.053 & 0.21877 & 0.21609 & 0.17038 \\
180 & $3.17\times10^{-6}$ & $3.26\times10^{-6}$ & 1.028 & 0.50000 & 0.50000 & 0.50000 \\
200 & $5.45\times10^{-6}$ & $5.22\times10^{-6}$ & 0.957 & 0.50000 & 0.50000 & 0.50000 \\
220 & $1.23\times10^{-5}$ & $1.21\times10^{-5}$ & 0.981 & 0.50000 & 0.50000 & 0.50000 \\
240 & $2.47\times10^{-5}$ & $2.32\times10^{-5}$ & 0.940 & 0.44359 & 0.50000 & 0.50000 \\
\bottomrule
\end{tabular}

\caption{Original and shifted-window MC results. Limits use the inherited
coupling coordinate. The last column is the Gaussian benchmark centered
at the same MC core. All probabilities are local asymptotic references.}
\end{table}
\begin{table}[H]\centering\small
\begin{tabular}{rrrrrr}
\toprule
$m_0$ & Gaussian at $c$ limit & MC/G at $c$ & $r_{MC}-r_G$ & Old MC fit (\%) & New MC fit (\%) \\
\midrule
60 & $4.54\times10^{-6}$ & 2.548 & -0.520 & 39.26 & 39.65 \\
80 & $6.20\times10^{-6}$ & 1.232 & -0.284 & 67.56 & 70.95 \\
100 & $2.19\times10^{-6}$ & 1.052 & -0.292 & 73.19 & 78.21 \\
120 & $2.12\times10^{-6}$ & 1.088 & -0.076 & 74.60 & 80.37 \\
140 & $1.84\times10^{-6}$ & 1.136 & +0.068 & 76.52 & 81.30 \\
160 & $4.76\times10^{-6}$ & 1.059 & -0.167 & 75.22 & 81.75 \\
180 & $3.27\times10^{-6}$ & 0.997 & -0.434 & 75.27 & 80.46 \\
200 & $4.27\times10^{-6}$ & 1.221 & -0.002 & 75.71 & 79.37 \\
220 & $8.91\times10^{-6}$ & 1.355 & +0.073 & 73.86 & 77.30 \\
240 & $1.63\times10^{-5}$ & 1.427 & -0.048 & 72.38 & 74.97 \\
\bottomrule
\end{tabular}

\caption{Matched centered-window comparison and retained MC probability
in the fitted bins. Support normalization remains unchanged; shifting
the mask does not renormalize its retained template fraction to one.}
\end{table}

At the native 80~MeV hypothesis, the reconstructed center is 78.10~MeV.
The shifted MC upper limit is $7.639\times10^{-6}$, compared with
$6.362\times10^{-6}$ before shifting, a 20.08\% increase. The local
probability decreases from 0.02477 to 0.007875; an increased upper limit
and a smaller local probability are compatible because they answer
different likelihood questions.

On the dense exploratory MC grid the smallest local probability moves
to 50~MeV ($p_0=0.003798$). That template uses the previously failing
low-mass interpolation region. It is not promoted to a validated
search result; the native comparison above remains the primary result.
The unchanged selected-candidate association and v13/v16 production
qualifications still apply. No global probability, coverage calibration
or physically validated A-prime exclusion follows from this recentering.

\clearpage
\subsection{Reconstructed cores and window placement}
\fig{.76}{../figures/core_centered_windows.pdf}{Native selected-MC cores
and the original versus shifted $\pm2.25\sigma_{\rm nominal}$ windows.
Blue curves are unchanged empirical templates. The local Gaussian plus
affine-pedestal fit (red) estimates the center only and is never used to
replace the full MC signal shape in the observed likelihood. The red
vertical line marks the fitted core center.}
